\frontmatter % Use roman page numbering style (i, ii, iii, iv...) for the pre-content pages

\pagestyle{plain} % Default to the plain heading style until the thesis style is called for the body content

%----------------------------------------------------------------------------------------
%	TITLE PAGE
%----------------------------------------------------------------------------------------

\begin{titlepage}
\begin{center}

\vspace*{0.16\textheight}
%{\scshape\LARGE \univname\par}\vspace{1.5cm} % University name
% \bigskip  
%\textsc{\Large Doctoral Thesis}\\[0.5cm] % Thesis type

%\HRule \\[0.4cm] % Horizontal line

%{\huge \bfseries \ttitle\par}\vspace{0.6cm} % Thesis title automatic

{\bfseries
\LARGE{Seismic Event Coda-Correlation Imaging of the Earth's Interior}\\
%\bigskip
%\large{Astrophysical, Geochemical and Biological Constraints on Habitability}
}

\vspace{1cm} % Thesis title

%\HRule \\[1.2cm] % Horizontal line
 
\begin{minipage}[t]{0.3\textwidth}
\begin{center}
 \large
%\emph{Author:}\\
 \authorname % Author name - remove the \href bracket to remove the link
\end{center}
\end{minipage}

%\begin{minipage}[t]{0.4\textwidth}
%\begin{flushright} \large
%\emph{Supervisor:} \\
%\href{http://www.jamessmith.com}{\supname} % Supervisor name - remove the \href bracket to remove the link  
%\end{flushright}
%\end{minipage}\\[3cm]

\vspace{6.3cm}
%\vspace{1cm}
%\includegraphics[width=0.6\textwidth]{gfx/front_cover.pdf} \\ \bigskip  
%\vspace{1cm}

\large \textit{A thesis submitted for the degree of \degreename}\\[0.7cm] % University requirement text

\vspace{0.0cm}

%\textit{in the}\\[0.4cm]
\groupname

\deptname % Research group name and department name

\vspace{-0.3cm}

\univname % University name

%{\small Draft (\today)}\\[2cm] % Date
%\includegraphics{Logo} % University/department logo - uncomment to place it

\includegraphics[width=6cm]{gfx/ANU_LOGO_cmyk_56mm.pdf}

\vspace{-0.6cm}

{\normalsize March, 2022}\\[2cm] % Date

\end{center}
\end{titlepage}

%----------------------------------------------------------------------------------------
%	DECLARATION PAGE
%----------------------------------------------------------------------------------------

\begin{declaration}
%\addchaptertocentry{\authorshipname} % Add the declaration to the table of contents
\begingroup
\large
\noindent This thesis presents a collection of research works conducted between May 2018 and February 2022 at the Research School of Earth Sciences, The Australian National University, Canberra, Australia. Except where acknowledged in the text, the material presented in this thesis is, to the best of my knowledge, original and has not been submitted in whole or part for a degree in any university. Co-authors have given formal consent for all papers to be presented in this thesis.
%\begin{itemize}
%\item Chapter~\ref{ch:HabitableWorlds} is the paper \cite{Lineweaver2012AnnRev}. Section~\ref{sec:HZones} is my major research contribution to the paper. Research on the distribution of biomass on Earth, the constraints for habitable zones within the Earth, and the energy requirements and sources for the earliest bacteria is my work.
%\item Chapter~\ref{ch:biologicalhab} is the paper \cite{Chopra2016} where my research forms a major contribution to all sections of the paper.
%\end{itemize}


\bigskip
\vspace{1cm}
\begin{flushright}
\authorname\\
\today
\end{flushright}
\endgroup
\end{declaration}
\vspace{3cm}
\begingroup
\endgroup
\cleardoublepage

% \begin{declaration}
% \addchaptertocentry{\authorshipname} % Add the declaration to the table of contents
% \noindent I, \authorname, declare that this thesis titled, \enquote{\ttitle} and the work presented in it are my own. I confirm that:

% \begin{itemize} 
% \item This work was done wholly or mainly while in candidature for a research degree at this University.
% \item Where any part of this thesis has previously been submitted for a degree or any other qualification at this University or any other institution, this has been clearly stated.
% \item Where I have consulted the published work of others, this is always clearly attributed.
% \item Where I have quoted from the work of others, the source is always given. With the exception of such quotations, this thesis is entirely my own work.
% \item I have acknowledged all main sources of help.
% \item Where the thesis is based on work done by myself jointly with others, I have made clear exactly what was done by others and what I have contributed myself.\\
% \end{itemize}
 
% \noindent Signed:\\
% \rule[0.5em]{25em}{0.5pt} % This prints a line for the signature
 
% \noindent Date:\\
% \rule[0.5em]{25em}{0.5pt} % This prints a line to write the date
% \end{declaration}

% \cleardoublepage

%----------------------------------------------------------------------------------------
%	QUOTATION PAGE
%----------------------------------------------------------------------------------------

% \vspace*{0.2\textheight}

% \noindent\enquote{\itshape Thanks to my solid academic training, today I can write hundreds of words on virtually any topic without possessing a shred of information, which is how I got a good job in journalism.}

% \hfill Dave Barry

%----------------------------------------------------------------------------------------
%	ACKNOWLEDGEMENTS
%----------------------------------------------------------------------------------------

\begin{acknowledgements}
\addchaptertocentry{\acknowledgementname} % Add the acknowledgements to the table of contents
\vspace{0.4cm}
\begingroup
\normalsize
Many years later, as looking back, I will still remember the picture of a sunny and warm autumn afternoon when golden leaves fall quietly outside the window, and Hrvoje on the other side of an old desk asked me, ``\textit{...do you accept the challenge...}''. A challenging but fascinating journey with Prof. Hrvoje Tkalčić, my supervisor, began that afternoon. Hrvoje provides me with meticulous academic guidance and helps throughout my PhD progress for tackling the coda-correlation tomography problem. On the other hand, he always inspires me to ``\textit{think big}'', for so many times when I was immersed and deeply captured in a one-track mind. In this journey, I have learned a lot from him. The picture will be in my mind forever. Thank you, Hrvoje!

I am also grateful to Prof. Ian Jackson, Dr Caroline M. Eakin, and Dr Andrew Valentine in my supervisory panel. They always help me reflect on what I have been doing along with my PhD progress. Their insightful comments and suggestions are always beneficial and helpful for completing this thesis. Talking with them can be like finding simple and quick answers on Google, or more often like being immersed in the great ocean of knowledge. Thank you, my supervisory panelists!

I feel lucky and glad to be a part of the Geophysics group (the former Seismology and Mathematical Geophysics group) at the Research School of Earth Sciences (RSES). It is a great environment for learning, research, conversations or debates, and more. It is also a wonderful scene beyond the ivory tower, where I get trained, improved, and connected to the scientific community and real-world sciences. Moreover, the RSES is a warm place to meet friends to share joys and sorrows, with many pictures painted in my mind and never fade!

I am also grateful to the resources and technical support from the National Computational Infrastructure (NCI) and the Terrawulf facility at the RSES for completing this thesis. The Australian Government provides the NCI infrastructure through the National Collaborative Research Infrastructure Strategy scheme (NCRIS). The Terrawulf is developed with support from AuScope, Australia's lead provider of research infrastructure, and funded through the NCRIS.

Finally, I would like to thank my parents, sister, and grandparents. Across a broad ocean, their support and encouragement are always with me!
\endgroup
\end{acknowledgements}



%----------------------------------------------------------------------------------------
%	ABSTRACT PAGE
%----------------------------------------------------------------------------------------

\begin{abstract}
\addchaptertocentry{\abstractname} % Add the abstract to the table of contents
\begingroup
\normalsize
Seismic coda waves are the late part of the seismic energy generated by earthquakes. Global coda correlograms are constructed by cross-correlating and stacking seismic event late coda records that are noisy and seemingly useless, but they exhibit many prominent features sensitive to the Earth's internal structure. Thus, the coda correlation rises as a new paradigm in global observational seismology. As a new category of observation, the correlation features, if interpreted correctly, can provide new information about the Earth's interior.

How to accurately utilise seismic event coda correlations, for instance, in “global coda correlation tomography,” has been controversial and unresolved. Some attempts treat coda correlations as reconstructed seismic waves, which is on a par with methods developed in ambient-noise correlations, for they share similar data processing and computation routines. However, that introduces erroneous interpretation because theoretical analyses have demonstrated fundamental differences in the formation mechanisms of coda correlations and ambient-noise correlations. Therefore, we need a solution, a correct approach, to allow us to use a massive amount of coda correlation observables to increase constraints on the Earth's interior.

This thesis consists of theoretical analyses, method developments, and applications for utilising seismic event coda correlations to image the Earth's interior. We first conduct comprehensive analyses to quantitatively `dissect' coda correlations for their formation mechanism. The analyses reveal the mathematical relationship between coda correlations and the Earth's internal structure. Based on that, we build a novel framework toward global coda correlation tomography. We verify the new framework in experiments and compare it with the method based on the assumption of seismic wave reconstruction. We illustrate significant inaccuracy in tomographic images can arise if coda correlations are treated as reconstructed seismic waves. Then, in an application, we provide a new class of observation for inner-core shear-wave anisotropy utilizing coda correlations in the new framework. We find that inner-core shear waves travel faster by at least ~5 s in directions oblique to the Earth's rotation axis than directions parallel to the equatorial plane (anisotropy of >0.8\%). Our inner-core shear-wave anisotropy observations place new constraints on the inner core mineral composition. Finally, we extend the principles to cross-correlations between source events and devise a new way to build global inter-source correlations. We demonstrate that a single seismic station is sufficient to construct a global correlogram. The correlogram exhibits prominent features sensitive to the internal planetary structures. We show implementations to constrain the Earth's and Martian cores' sizes and confirm a large Martian core. This provides a new paradigm for imaging planetary interiors on a global scale with currently realizable resources in planetary missions.
\bigskip

%\textbf{Keywords:} \keywordnames

\endgroup
\end{abstract}

%----------------------------------------------------------------------------------------
%	LIST OF CONTENTS/FIGURES/TABLES/Equations PAGES
%----------------------------------------------------------------------------------------

\renewcommand{\listfigurename}{List of Figures}
\renewcommand{\listtablename}{List of Tables}

\begin{spacing}{0.94} 
\tableofcontents % Prints the main table of contents
\end{spacing}

\listoffigures % Prints the list of figures
\listoftables % Prints the list of tables
% \listoflistings % Prints the list of python codes

% Sheng
%\makeatletter
\listofmyequations %Write out the List of Equations
\addcontentsline{toc}{chapter}{\listequationsname}
%\makeatother

%----------------------------------------------------------------------------------------
%	ABBREVIATIONS
%----------------------------------------------------------------------------------------

%\begin{abbreviations}{ll} % Include a list of abbreviations (a table of two columns)
%\textbf{ANU}&\textbf{A}ustralian \textbf{N}ational \textbf{U}niversity\\
%\textbf{RSES}&\textbf{R}esearch \textbf{S}chool of \textbf{E}arth \textbf{S}ciences\\
%\end{abbreviations}

%%----------------------------------------------------------------------------------------
%%	PHYSICAL CONSTANTS/OTHER DEFINITIONS
%%----------------------------------------------------------------------------------------
%
%\begin{constants}{lr@{${}={}$}l} % The list of physical constants is a three column table
%
%% The \SI{}{} command is provided by the siunitx package, see its documentation for instructions on how to use it
%
%Speed of Light & $c_{0}$ & \SI{2.99792458e8}{\meter\per\second} (exact)\\
%%Constant Name & $Symbol$ & $Constant Value$ with units\\
%
%\end{constants}

%%----------------------------------------------------------------------------------------
%%	SYMBOLS
%%----------------------------------------------------------------------------------------
%
%\begin{symbols}{lll} % Include a list of Symbols (a three column table)
%
%$a$ & distance & \si{\meter} \\
%$P$ & power & \si{\watt} (\si{\joule\per\second}) \\
%%Symbol & Name & Unit \\
%
%\addlinespace % Gap to separate the Roman symbols from the Greek
%
%$\omega$ & angular frequency & \si{\radian} \\
%
%\end{symbols}

%----------------------------------------------------------------------------------------
%	DEDICATION
%----------------------------------------------------------------------------------------

%\dedicatory{For/Dedicated to/To my\ldots} 


